\section{Before anything else}
\label{sec:preface}

\subsection{Who this is for}

This note is written for someone who ships software. It assumes fluency with
SQL, comfort with at least one of LINQ, XQuery, list comprehensions or a
collection pipeline API, and the ordinary working knowledge of concurrency that
comes from having debugged it. It assumes no category theory, no process
algebra, and no prior contact with the rho calculus.

There is a companion document, \emph{Graph-Structured Lambda Theories: A
Reference for the Working Developer}~\cite{GSLTintro}, which covers overlapping ground for a
reader who wants the theory first and is willing to spend a chapter on
bisimulation to get it. This note is the other entrance to the same building. A
reader who finishes here and wants proofs will find them there; a reader who
starts there and wants to write code will find that here. Neither is a
prerequisite for the other.

\subsection{Why the order is backwards}

The convention is to teach a language's data definition language first --- here
are the types, here is how to make values --- and its control flow language
second. The convention is defensible and this note ignores it.

The reason is that the control flow language of rholang has a familiar door and
the data definition language does not. Every working developer has written a
query comprehension. Walking through that door and following the shape until it
breaks generates, in order, exactly the questions the data definition language
answers. Taught the other way round, the \texttt{language!} block looks like an
unusually elaborate schema declaration and its point is invisible.

So the order here is: the comprehension; what a comprehension cannot do; what
rholang adds; what that addition forces us to have; and then, twice, what
happens when we widen the thing that was added.

\subsection{The running example}

One example runs from beginning to end. A \emph{forager} moves through an
environment of specimens. Some specimens are rich and worth the cost of opening;
some are decoys and are not. The forager has a cheap structural test it can
apply, and that test is evidence rather than proof. It has a finite budget, it
spends from that budget to look and to open, and if the budget reaches zero the
forager stops --- permanently.

The example is drawn from a research programme on learning as computation, in
which such a forager is a very small instance of a \emph{mortal
scientist}~\cite{ScientistNote}: a
computation that forms hypotheses about its environment, pays to test them, and
is punished by the environment when they are wrong. The point of choosing it is
not the biology. It is that the same four lines of code are legible at every
stage of this note --- as a query, as a transaction, as a hypothesis, and
finally as a learning rule --- and that the numbers quoted in the last part were
measured by running it.

\subsection{Notation, and two conventions worth flagging}

The calculus underneath rholang is the \rhoc{} calculus~\cite{MR05}. The name abbreviates
\textbf{r}eflective \textbf{h}igher \textbf{o}rder calculus, and dropping the
periods gives the transliteration of the Greek letter after $\pi$ --- a joke
about what comes after the $\pi$-calculus. It is written ``\rhoc{} calculus'',
never with the Greek letter.

Quotation is written \texttt{@P}: the name that stands for the process
\texttt{P}. Dereference, sometimes called dropping or unquoting, is written
\texttt{*x}: the process that the name \texttt{x} stands for. These are inverse
in one direction on the nose --- \texttt{@*x} is \texttt{x} --- and that fact
does more work later than it looks like it should.

Code listings in this note come from two places, and it is worth knowing which
is which. Listings labelled as running examples are taken verbatim from the
\texttt{f1r3node-rust} repository~\cite{f1r3node}, where they compile and are covered by tests.
Listings that illustrate a design under development are marked as such. Nothing
here silently pretends that a plan is a product; \S\ref{sec:status} is a table
of what actually runs.
